% !TeX program = xelatex
% =============================================================
% 模板 08：中文本科毕业论文（ctexbook + XeLaTeX + BibTeX / GB/T 7714）
% 编译： scripts\build.ps1 cn-thesis
% 结构： 封面 → 中英文摘要 → 目录 → 五章正文 → 参考文献 → 致谢 → 附录
% 版式： A4，上下 2.54cm、左右 3.0cm，正文小四宋体 1.5 倍行距
% 编号： 图/表/公式按章编号（图 1-1、表 2-1、式 (3-1)）
% =============================================================
\documentclass[UTF8,a4paper,zihao=-4,fontset=windows,openany]{ctexbook}

\usepackage{geometry}
\geometry{top=2.54cm,bottom=2.54cm,left=3.0cm,right=3.0cm}
\usepackage{amsmath,amssymb,bm}
\usepackage{booktabs}
\usepackage{array}
\usepackage{multirow}
\usepackage{graphicx}
\usepackage{caption}
\usepackage{fancyhdr}
\usepackage{gbt7714}
\usepackage[hidelinks]{hyperref}

\linespread{1.5}

% 页脚居中页码，不要页眉
\pagestyle{fancy}
\fancyhf{}
\fancyfoot[C]{\thepage}
\renewcommand{\headrulewidth}{0pt}

\captionsetup{font=small,labelsep=quad}
\captionsetup[table]{position=top}
\captionsetup[figure]{position=bottom}

% ------------------------- 图表公式按章编号 -------------------------
% 想要全文连续编号，删掉下面 6 行即可（编号变成 图1、表1、式(1)）
\counterwithin{figure}{chapter}
\counterwithin{table}{chapter}
\counterwithin{equation}{chapter}
\renewcommand{\thefigure}{\thechapter-\arabic{figure}}
\renewcommand{\thetable}{\thechapter-\arabic{table}}
\renewcommand{\theequation}{\thechapter-\arabic{equation}}

% 参考文献标题由 natbib 的 \bibsection 生成，在这里补目录项，页码才不会错位
\renewcommand{\bibsection}{%
  \chapter*{\bibname}%
  \phantomsection
  \addcontentsline{toc}{chapter}{\bibname}%
}

% ------------------------- 章节标题样式 -------------------------
\ctexset{
  chapter = {
    format    = \zihao{3}\heiti\centering,
    name      = {第,章},
    number    = \arabic{chapter},
    aftername = \quad
  },
  section    = {format = \zihao{4}\heiti},
  subsection = {format = \zihao{-4}\heiti}
}

\begin{document}

% ============================ 封面 ============================
% 不用 titlepage 环境：它自带 \newpage + 重置页码，会和后续章节的空页冲突
\thispagestyle{empty}
\begin{center}
  \vspace*{1.8cm}
  {\zihao{2}\heiti XX 大学\par}
  \vspace{0.6cm}
  {\zihao{3}\heiti 本科毕业论文（设计）\par}
  \vspace{2.8cm}
  {\zihao{2}\heiti 占位题目：基于 XXX 方法的 YYY 问题研究\par}
  \vspace{2.8cm}
  {\zihao{4}
  \renewcommand{\arraystretch}{1.6}
  \begin{tabular}{rl}
    学\qquad 院： & XX 学院 \\
    专\qquad 业： & XX 专业 \\
    班\qquad 级： & 2026 级 X 班 \\
    学生姓名： & 张\quad 三 \\
    学\qquad 号： & 2026000000 \\
    指导教师： & 李\quad 四\quad 副教授 \\
  \end{tabular}\par}
  \vfill
  {\zihao{4} 2026 年 6 月\par}
\end{center}

% ============================ 前置部分 ============================
% 封面计为第 0 页且不显示页码，摘要页从 I 开始；不要用 \frontmatter，会和封面抢分页
\pagenumbering{Roman}
\setcounter{page}{0}

% ------------------------- 中文摘要 -------------------------
\chapter*{摘\quad 要}
\addcontentsline{toc}{chapter}{摘要}

这里是中文摘要占位文字。本文围绕 XXX 问题，采用 YYY 方法建立了 ZZZ 模型，并通过实验验证了模型的有效性。本科毕业论文摘要一般 300--500 字，应能独立成篇，不出现图表和公式编号，不写「本文第一章」之类的导航语。

\vspace{1em}
\noindent{\heiti 关键词：}占位关键词一；占位关键词二；占位关键词三；占位关键词四

% ------------------------- 英文摘要 -------------------------
\chapter*{Abstract}
\addcontentsline{toc}{chapter}{Abstract}

This is a placeholder abstract. This thesis studies the XXX problem with the YYY method, builds a ZZZ model, and validates it experimentally. Replace this paragraph with your own 200--300 word abstract, which should be self-contained and contain no figure or equation numbers.

\vspace{1em}
\noindent{\bfseries Keywords:} placeholder keyword one; placeholder keyword two; placeholder keyword three

% ------------------------- 目录 -------------------------
\tableofcontents

% ============================ 正文 ============================
\mainmatter   % 阿拉伯数字页码，从 1 开始

\chapter{绪论}
\label{chap:intro}

\section{研究背景与意义}

这里是研究背景占位内容。说明课题来自什么工程/学术问题、为什么值得做。引用文献示例：文献~\cite{zhang2024demo} 给出了相关工作，英文文献见~\cite{smith2023demo}。

\section{国内外研究现状}

按「国外—国内—小结」分段写，每段结尾点明已有工作的不足，自然过渡到本文工作。

\section{本文主要工作}

本文主要工作包括：
\begin{enumerate}
  \item 建立了 XXX 模型；
  \item 设计了 YYY 方法；
  \item 通过实验验证了上述方法的有效性。
\end{enumerate}

\section{论文结构安排}

第~\ref{chap:intro} 章为绪论；第~\ref{chap:method} 章介绍方法与模型；第~\ref{chap:design} 章给出设计实现；第~\ref{chap:experiment} 章为实验与结果分析；第~\ref{chap:conclusion} 章总结全文。

% ============================ 第 2 章 ============================
\chapter{相关理论与方法}
\label{chap:method}

\section{问题描述}

设输入变量为 $x_i$，权重为 $\alpha_i$，则加权综合指标可写为
\begin{equation}
  S(x)=\sum_{i=1}^{n}\alpha_i x_i+\beta_0,
  \label{eq:weighted}
\end{equation}
其中 $\beta_0$ 为偏置项。式~\eqref{eq:weighted} 为公式占位示例。

\section{模型建立}

实验参数见表~\ref{tab:params}。

% ------------------------- 表格 -------------------------
\begin{table}[htbp]
  \centering
  \caption{实验参数设置（表格占位示例）}
  \label{tab:params}
  \begin{tabular}{lccc}
    \toprule
    参数 & 符号 & 取值 & 单位 \\
    \midrule
    学习率     & $\eta$    & 0.01 & -- \\
    迭代次数   & $N$       & 1000 & 次 \\
    正则系数   & $\lambda$ & 0.10 & -- \\
    \bottomrule
  \end{tabular}
\end{table}

\section{本章小结}

本章介绍了 XXX 理论与 YYY 方法，为第~\ref{chap:design} 章的设计提供依据。

% ============================ 第 3 章 ============================
\chapter{系统设计与实现}
\label{chap:design}

\section{总体方案}

系统总体框图见图~\ref{fig:framework}。

% ------------------------- 图 -------------------------
% 换成真实图片：把 \rule 那行替换为
% \includegraphics[width=0.75\textwidth]{figs/framework.pdf}
\begin{figure}[htbp]
  \centering
  \framebox[0.65\textwidth][c]{\rule{0pt}{4cm}}
  \caption{系统总体框图（占位图，请替换为真实图片）}
  \label{fig:framework}
\end{figure}

\section{关键模块实现}

分模块说明输入、处理、输出，配合伪代码或流程说明。

\section{本章小结}

本章完成了系统设计，第~\ref{chap:experiment} 章给出实验验证。

% ============================ 第 4 章 ============================
\chapter{实验与结果分析}
\label{chap:experiment}

\section{实验设置}

说明数据集、软硬件环境、对比方法与评价指标。参数配置见第~\ref{chap:method} 章表~\ref{tab:params}。

\section{结果分析}

结合图表说明：主要结果是什么、与对比方法相比好在哪里、误差来源是什么、结果是否稳定。

\section{本章小结}

本章实验表明 XXX 方法在 YYY 指标上优于对比方案。

% ============================ 第 5 章 ============================
\chapter{总结与展望}
\label{chap:conclusion}

\section{全文总结}

概括本文完成的三项工作和主要结论。

\section{不足与展望}

指出方法的局限（数据规模、应用场景、计算开销等），给出可行的改进方向。

% ============================ 参考文献 ============================
\bibliographystyle{gbt7714-numerical}
\bibliography{refs}

% ============================ 致谢 ============================
\chapter*{致\quad 谢}
\addcontentsline{toc}{chapter}{致谢}

感谢指导教师在选题、方案与写作上的指导，感谢实验室同学在实验与数据处理上的帮助，感谢家人一直以来的支持。

% ============================ 附录 ============================
\appendix
\ctexset{chapter={name={附录,},number=\Alph{chapter}}}
\chapter{核心代码}
\label{app:code}

附录放不宜进正文的材料：核心代码、原始数据表、详细推导等。

\end{document}
